\documentclass[11pt]{article}

\usepackage[margin=1in]{geometry}
\usepackage{amsmath, booktabs, graphicx, array, tabularx, float, placeins}
\usepackage[hidelinks]{hyperref}
\usepackage{xurl}

\title{Corpses and Corpora: Can an AI Learn From a Graveyard\\
of Failed Basketball Studies?}
\author{Rami Zheman \\ Independent Researcher, Alexandria, VA}
\date{November 2026}

\begin{document}
\maketitle

\begin{abstract}
Many of my basketball research ideas failed to produce the findings I expected, and I kept the records. This paper tests whether those records help an AI agent decide if a new research idea is worth running. Five language models judged ten of my past ideas, as of the day each was proposed, with no archive, with document search over 319 dated earlier studies, or with graph search. A good call meant stopping ideas that later proved untestable or an artifact of a design flaw, and letting the rest go ahead. Sports organizations are building systems of this kind, so the question matters for how they design their archives.

The archive did not reliably improve decisions: it fixed 47 wrong calls and broke 46 right ones. The models repeated earlier studies accurately but could not tell when a study applied to a new idea; one data-quality record bore on two ideas with opposite correct answers, and nothing in it said which one it invalidated. Better search did not help. Embeddings found the closest earlier study in about 40 to 50 percent of runs, against 2 of 45 for keyword search, and accuracy did not improve. Handing the models the relevant studies directly raised accuracy for four of five models, but ten ideas are too few to say whether that gain is real. A test with 60 to 100 ideas is the next step.
\end{abstract}

\section{Introduction}
\label{sec:intro}

This paper began as a failure. I spent months on one basketball topic, in-game adaptation, believed I had a strong finding, and wrote a full paper. Then a full start-to-finish review showed that the data I had used to answer the question was biased: I had pulled data that was narrowed by which results were available, and I had not noticed. I used AI to pursue many other questions, hoping to salvage the work. Over time, those questions and my other projects built up into an archive of 319 study families, each a record of what I tried and what happened. It covers many topics, of which adaptation is only one. A good share of it is a graveyard of ideas that did not work, though not all of it: some studies held up.

That archive suggested a question. If the corpses can be read, can they help someone avoid the next corpse? Put more formally, if an AI agent reads the record of past studies, does it make a better decision about whether a new idea is worth running? The corpses become the corpus, and the paper asks whether the corpus is any use. Sports organizations are starting to build systems that accumulate knowledge in this way (Buffalo Sabres 2026; Sports Business Journal 2025), which is part of why the question matters, but the question itself came from my own archive.

The short answer is that models could retrieve and accurately repeat my past studies, but they did not reliably tell when a past study applied to a new idea. Across five models and ten tasks, the archive fixed 47 wrong decisions and broke 46 right ones. Search that returned the closest earlier study far more often (18 and 22 of 45 runs, against 2) did not improve accuracy, and handing over the relevant studies pointed upward (Section~\ref{sec:study1b}). In the clearest case, one data-quality record bore on two ideas with opposite correct answers, and nothing in it said which one it invalidated. That points to a concrete design change: records need to say when their lessons apply. The setup is one program's full archive, failures included, with each task dated to the day the idea was proposed and graded against what happened when it was run.

I found no earlier study that tests this on an organization's own archive, failures included.

\subsection*{What the agent is asked to do}

Each task is a research idea as of the day I proposed it. I know how it turned out because I later ran it. The model must decide, as of that day, whether to go ahead or stop, and it can search only studies dated before that day; the task's own study and anything too close to its results are removed. Its call is graded against what happened when I ran the idea, which I call the \emph{answer key} (Section~\ref{sec:tasks}). With no memory the search tool returns nothing, so the model decides from the idea alone; every model is also run that way as the baseline.

One task shows the judgment. The pick-and-roll drive study (RT-012) asks whether a teammate's drive just before a pick-and-roll makes that pick-and-roll score better, and what the model could retrieve was an earlier study of how possessions are classified by play type. The model must judge whether that labeling problem affects the measurements the new study needs. When I later ran it, the apparent gain disappeared under a different classification, so the answer key says stop (Section~\ref{sec:rt012}).

A second task has the other answer. The lineup shot-mix study (RT-112) averages across substitutions, and no players are named. For each substitution it compares the share of shots that are three-pointers for the players who come in with that share for the players they replace, then subtracts the change at stoppages where the same players stay. When I ran it, the incoming group averaged $+3.1$ percentage points more of its shots as threes than that comparison. The same measurement did not move when only the opponent substituted, so the design can see a change in personnel and the answer key says go ahead. The archive had no earlier study of that same question for the model to retrieve.

The paper tests one step: whether reading the record of prior studies improves a single go-ahead-or-stop decision. Ingestion, continuous operation, and write-back are left to Section~\ref{sec:future}. Study 1 compares no memory with document search and graph search, using either keyword matching or embeddings, and Study 1b hands the models the relevant studies directly, so search is removed as the weak point.

\section{Related work}
\label{sec:related}

\subsection{Sports organizations}

Sports organizations hold a large stock of analytical capital (game logs, tracking, video, scouting observations, and the record of what was tried, what worked, and what failed), and most of it sits in files, in people's heads, or in analyses left unused the next time a similar question comes up. Erhardt and Martin-Rios (2016) find that football organizations tend toward explicit, codified knowledge and ice hockey organizations toward tacit transfer through people, and codified records are the ones a system can retrieve. A survey of football analytics finds thin dedicated staffing (Science and Medicine in Football 2024). Public examples range from cloud storage (Mission Cloud 2024) to semantic search (Sports Business Journal 2025) to a platform meant to learn (Buffalo Sabres 2026). These reports leave untested whether retrieved history improves the next analytical decision. That is the gap this paper measures, on a basketball research-design task rather than in a live front office.

\subsection{Case-based reasoning}

Case-based reasoning solves a new problem by retrieving a similar past case and reusing its solution (Aamodt and Plaza 1994). Its standard cycle has four steps: retrieve, reuse, revise, and retain. The revise step tests whether the reused solution actually fits the new problem before it is kept. The benchmark here measures retrieve and reuse with no separate revise step: the same agent that retrieves a prior study decides whether it applies. The cases in Section~\ref{sec:transfer} are examples of what can happen without that step.

\subsection{Retrieval-augmented generation and irrelevant context}

Lewis et al.\ (2020) augment a language model with retrieved documents so that answers can use material outside the model's parameters. The usual target is factual grounding. Research memory asks whether the quoted study applies. A model can quote a prior study accurately and still apply it to a question outside that study. The benchmark separates the two: citation checks ask whether the record was reproduced, and the decision score asks whether the decision was right.

Shi et al.\ (2023) show that adding irrelevant information to a problem statement degrades language-model reasoning, even when the relevant information is unchanged. Retrieved research records are a natural source of such context: a record can be on topic and still have no bearing on the decision. The decision shifts reported here are consistent with that finding, although this benchmark does not isolate distraction from other causes.

SoundnessBench reconstructs 1,099 machine-learning research proposals from ICLR submissions and tests 12 frontier models on judging soundness before execution (Ho et al.\ 2026). The models tend to approve flawed proposals. More aggressive prompting mostly moves the errors from false approvals to false rejections. The benchmark judges each proposal alone, with no retrieved history. ForeSci has 500 tasks in four AI domains, and each task comes with a knowledge base limited to evidence from before a cutoff date (Tian et al.\ 2026). It reports evidence-decision decoupling: an agent can cite relevant evidence and still predict the wrong research object. The caution shift in Section~\ref{sec:secondary} resembles the trade-off SoundnessBench reports. This paper's finding that models cite a record accurately but do not reliably judge whether it applies is consistent with ForeSci's decoupling. This paper differs in using one organization's own archive with recorded outcomes.

\subsection{Memory of failures in language agents}

Reflexion (Shinn et al.\ 2023) has an agent write verbal reflections on failed attempts into an episodic memory and reuse them on later attempts, improving performance on decision-making, reasoning, and coding tasks. The lessons there are reused on retries of closely related problems. The archive here holds lessons from different studies, and the question is whether a lesson recorded for one study transfers to another. A lesson such as ``this study was inconclusive because the play-type labels were unreliable'' is useful only if the next decision checks whether the new study depends on those labels in the same way.

\section{The archive and the benchmark}
\label{sec:archive}

\subsection{What the archive contains}

The archive is my own past work and papers: 319 studies from a program of basketball analytics covering four NBA seasons (2022--23 through 2025--26), regular season and playoffs. The studies cover ten topics, and Table~\ref{tab:archive_topics} shows how many there are in each and how each one turned out.

\begin{table}[ht]
\centering
\caption{Archive composition by topic area}
\label{tab:archive_topics}
\small
\begin{tabular}{lrrrrrr}
\toprule
Topic area & Fam. & Sup. & Rej. & Untest. & Art. & Uncert. \\
\midrule
Defensive suppression and counter strategy & 62 & 13 & 10 & 9 & 8 & 8 \\
Player absence and offensive reorganization & 42 & 10 & 12 & 5 & 3 & 7 \\
Structural roles and influence & 41 & 11 & 7 & 8 & 4 & 5 \\
Lineup and matchup effects & 30 & 3 & 7 & 7 & 3 & 6 \\
Play-type labeling and data quality & 18 & 1 & 1 & 2 & 6 & 2 \\
NCAA tournament prediction & 14 & 6 & 1 & 1 & 2 & 3 \\
In-game and series adaptation & 12 & 0 & 2 & 0 & 3 & 5 \\
Possession sequences and transitions & 10 & 2 & 3 & 1 & 3 & 0 \\
Playing time and luck response & 3 & 2 & 0 & 0 & 1 & 0 \\
Other (including internal evaluation) & 87 & 9 & 14 & 16 & 11 & 16 \\
\midrule
Total & 319 & 57 & 57 & 49 & 44 & 52 \\
\bottomrule
\end{tabular}

\vspace{0.4em}
\begin{minipage}{0.92\linewidth}\footnotesize
Sup.\ = supported, Rej.\ = rejected (adequately powered null), Untest.\ = untestable (data cannot support a test), Art.\ = artifact (apparent result produced by a design flaw), Uncert.\ = uncertain (inconclusive or inadequately tested). An additional 60 families are classified as not a test or no result recorded and are not in the outcome groups above.
\end{minipage}
\end{table}

It is a mixed body of work. Of the 259 studies with a recorded outcome, 57 held up and 57 were rejected by an adequately powered test. The other 145 were an artifact, untestable, or inconclusive.

\subsection{What a study family record contains}

Each study family is a structured record. It carries the research question and hypothesis, the dataset and sample size, the method and outcome measures, the primary result, an outcome status, the failure mechanism where one was identified, confounders and limitations, a validation status, related families, and source references.

The same facts are rendered twice, once as a document for search and once as a node in a graph (Appendix~\ref{app:history}). The outcome statuses are rule-based classifications of the records, and the benchmark measures agreement with that classification.

\subsection{The ten tasks}
\label{sec:tasks}

Each task presents a proposed basketball study. The agent sees the research question, the data available on the task date, and the data sources, and must decide whether the design should go ahead or stop. The answer key says stop when the task's own study was recorded as untestable or an artifact; supported, rejected, and uncertain studies are keyed as go ahead, because a test of the question as posed was runnable and informative.

Four tasks should be stopped and six should go ahead. The ten were selected by a predefined script from 196 eligible study families, and the list was fixed before the primary experiment (Appendix~\ref{app:selection}).

Table~\ref{tab:tasks_plain} describes each task in one plain line. Appendix Table~\ref{tab:tasks} gives the research questions as posed and the recorded result behind each decision. The models did not see the recorded result.

\begin{table}[ht]
\centering
\caption{The ten tasks in plain language. ``Stop'' means the answer key says the idea should be stopped.}
\label{tab:tasks_plain}
\small
\begin{tabularx}{\linewidth}{lXl}
\toprule
Task & What the idea asks, in plain words & Answer key \\
\midrule
RT-012 & Does a teammate's drive just before a pick-and-roll make that pick-and-roll score better? & Stop \\
RT-014 & Would following the system's recommended mix of play types win more games? & Stop \\
RT-065 & Can you tell in advance which teammate will take over a missing player's role? & Go ahead \\
RT-066 & Does the amount a team reshuffles its offense while a player is out predict how its scoring changes? & Go ahead \\
RT-101 & Do teams differ in how well their lineups cover offensive roles, apart from raw talent? & Go ahead \\
RT-112 & Does changing the five-man lineup change what kinds of shots a team takes? & Go ahead \\
RT-146 & After a defense shuts down pick-and-roll, does pull-up shooting reliably work as a counter? & Go ahead \\
RT-147 & Are the effects on an opponent's next possession large and stable enough to justify telling a team to avoid certain plays? & Go ahead \\
RT-179 & Is a team's ability to hold down opponents' second-half improvement a stable, measurable trait? & Stop \\
RT-183 & Do role players who mostly cut keep their role better after changing teams than players who mostly drive? & Stop \\
\bottomrule
\end{tabularx}

\end{table}



A relevant study shares the task's question or a recorded risk, by a fixed rule. Three tasks (the pick-and-roll drive study (RT-012), the opponent next-possession study (RT-147), the cutting-and-driving role study (RT-183)) have a close earlier study of the same question in memory; four (the play-mix recommendation study (RT-014), the offense-reshuffle study (RT-066), the lineup shot-mix study (RT-112), the second-half improvement study (RT-179)) have only earlier failures that share a recorded risk; and three (the missing-player role study (RT-065), the lineup role-coverage study (RT-101), the pull-up counter study (RT-146)) have nothing relevant (Appendix~\ref{app:retrieval}).

\section{Experimental design}
\label{sec:design}

\subsection{How the models received historical research}

I tested whether giving an AI model access to earlier research would help it decide whether a new research idea was worth pursuing. The same five models ran under each condition. I compared three basic approaches.

In the first, \emph{no memory}, the model received the new research question and a description of the available data. It had no access to the archive of earlier studies.

In the second, \emph{automatic retrieval}, the model could search the archive before making its decision. I tested two ways of organizing that search. Document search finds study records that match the search. Graph search can use links between concepts, hypotheses, and study records to find related studies beyond shared words. Within each, I also tested two ways of matching the query to the archive. Keyword matching looked for shared words. Embedding matching turned the query and each study into a list of numbers and ranked studies by how close those lists were. Keyword search returned up to four records. Embedding search returned up to ten. In both search setups the model first read the new proposal, wrote a query, and called the search tool. It then received the records found and made its recommendation. The prompt separates retrieval from the decision.

In the third approach, \emph{supplied studies}, the model was given the relevant earlier studies directly instead of having to find them (Section~\ref{sec:study1b}).

In every condition the model had the same job: assess the new research proposal and decide whether the idea should go ahead or stop. The difference was what historical evidence, if any, the model could use before that decision. Five models each judged the ten ideas three times (450 primary runs). The embedding search added 300 runs and did not repeat the no-memory condition. How each search works, the pilots, and the protocol checks are in Appendix~\ref{app:history}.

\subsection{Memory snapshots and the exclusion of each task's own result}

Each task's runs searched a per-task snapshot of the archive, built before any run. It contains only records dated before the task date and excludes the task's own study, records that cite it, probable duplicates, and records sharing distinctive result numbers with it (Appendix~\ref{app:snapshots}). Earlier studies of the same question that did predate the task were kept, because earlier work on the same question is legitimate evidence.

\subsection{How the decisions were scored}
\label{sec:outcome}

Each model had to decide whether a research idea should go ahead or stop. The benchmark compared that decision with the answer key (Section~\ref{sec:tasks}).

The benchmark scored the written recommendation together with the model's estimates of whether the study could be tested and whether its result was likely to be an artifact. A model could therefore recommend proceeding with changes but still receive a stop score if its estimates indicated that the study was unlikely to produce a reliable result. The exact scoring rules are in Appendix~\ref{app:scoring}.

I measured performance by averaging two results: how often the models correctly identified the four ideas that should stop, and how often they correctly identified the six that should go ahead. This prevented the larger group of go-ahead tasks from dominating the score. This measure, known as balanced accuracy, gives equal weight to both types of decision. After this section I call it the accuracy score.

Checking whether the conclusion changes under different scoring rules, and the exact statistical procedures, are in Appendix~\ref{app:scoring}.

\section{Primary results}
\label{sec:results}

Giving the models access to the archive produced no detectable improvement. None of the apparent improvements was strong enough to rule out random variation. Table~\ref{tab:pooled} gives the average accuracy score across the five models. Each row is one way of giving those models access to earlier research.

\begin{table}[ht]
\centering
\caption{Mean accuracy score across the five models on the ten research ideas. The score gives equal weight to stop decisions and go-ahead decisions. Descriptive; no uncertainty range.}
\label{tab:pooled}
\small
\begin{tabularx}{\linewidth}{l r X}
\toprule
Condition & Score & What the model could use \\
\midrule
No memory & 0.611 & Decides from the new proposal alone, with no past studies. \\
Automatic document retrieval (keyword) & 0.625 & Finds up to four records by matching words. \\
Automatic graph retrieval (keyword) & 0.639 & Finds up to four records through links between concepts, hypotheses, and studies. \\
Automatic document retrieval (embedding) & 0.558 & Finds up to ten records by how close their numerical representations are. \\
Automatic graph retrieval (embedding) & 0.578 & The same embedding ranking, through the graph, up to ten records. \\
Supplied studies & 0.725 & Receives the relevant earlier studies directly. \\
\bottomrule
\end{tabularx}
\end{table}

The score runs from 0 to 1. A 1 would mean the model matched the answer key on the ideas that should stop and on the ideas that should go ahead. A 0.5 is what guessing produces, because a guess is right about half the time on each kind of idea and the score averages those two rates. A 0.6 is only a little better than guessing. The 0.611 with no memory is that kind of result. The highest score is the 0.725, where the relevant studies were handed over, and it is still well short of getting the decisions right. Even when embedding search found more relevant earlier studies, the accuracy score did not rise.

Memory helped some models and hurt others (per-model values in Appendix~\ref{app:permodel}). There was no consistent improvement across all five models. For keyword search, the change from no memory ranged from a loss of 0.083 to a gain of 0.208. DeepSeek showed the largest improvement with graph search. The uncertainty range includes no difference. This result needs more tasks to confirm. Sonnet, the strongest model without memory, is lower with document search and about unchanged with graph search. None of these differences was statistically conclusive under any of the scoring rules. Results depend on how unreadable outputs are scored, most for Sonnet; those outputs failed the required format (Appendix~\ref{app:scoring}). The two models that reason before answering showed no consistent direction either (Appendix~\ref{app:more}).

\subsection{What the models did with the record}
\label{sec:secondary}

When citing earlier studies, the models generally reproduced their recorded outcome statuses accurately, but they sometimes claimed that earlier work had settled more than it had, and retrieved failures made them more cautious in both directions. No final answer cited an identifier the tool had not returned, and every status checked against the record matched (per-model counts in Appendix~\ref{app:more}). Some runs incorrectly claimed that an earlier study had already answered the question, on a task where no such study was in memory (eight tasks, 24 eligible runs per model and condition). That happened in 25 to 50 percent of eligible runs under document search and 42 to 67 percent under graph search.

Retrieved history also made models more cautious in both directions. In Experiment 1, document search raised Gemini's stop rate from 0.50 to 0.75 where stopping was correct, and from 0.22 to 0.39 where it was not (Appendix Table~\ref{tab:caution}). The gains and losses are similar in size, which is consistent with the accuracy score barely moving. The models' own estimates of whether a study could be tested ranked the tasks better with memory for four of the five models, but that ranking did not produce better stop-or-go decisions at the 50\% cutoff.

\subsection{Whether the search returned the earlier study}
\label{sec:embed}

The decision scores above mix a failed search with a failed judgment. A model can score badly because the search never returned the earlier study, or because the study came back and the model still made the wrong call. This section separates those two. It asks whether the search returned the earlier study. That score is already in Table~\ref{tab:pooled}.

Table~\ref{tab:embed_retr} covers the four search rows in Table~\ref{tab:pooled}: document search and graph search, each with keyword matching and with embeddings. The model wrote a query without being shown which studies a fixed rule already marked as relevant. I counted that afterward.

Only three of the ten ideas have one closest earlier study, a prior study of the same question. Keyword search almost never found the closest study; embedding search found it about 40 to 50 percent of the time (Table~\ref{tab:embed_retr}).

Seven of the ten ideas have at least one relevant earlier study, either that closest study or an earlier failure that shares a recorded risk. The other three ideas have nothing relevant in the archive, so the table covers the seven.

Those embedding conditions are the 0.558 and 0.578 in Table~\ref{tab:pooled}, below the 0.611 with no memory. The search returned the earlier study more often. The decision was not better. Those rows also change how many records come back, four against ten, and how much text each result can contain (Appendix~\ref{app:embed}). The scores combine the matching method with both of those limits. Returning ten studies instead of four may have added distraction, and these results cannot separate that from the question of whether the lesson applies. The per-model values and the limits are in Appendix~\ref{app:more}.

\begin{table}[H]
\centering
\caption{Whether the search returned an earlier study, for the four search rows in Table~\ref{tab:pooled}. The closest study is a prior study of the same question: 3 ideas, 5 models, 3 repeats, 45 runs. A relevant study shares the question or a recorded risk: 7 ideas, 105 runs. The numerator is runs in which the tool returned that study. Returned and cited is runs, out of 105, in which that study was both returned and cited.}
\label{tab:embed_retr}
\label{tab:embed_retr_full}
\footnotesize
\setlength{\tabcolsep}{4pt}
\begin{tabularx}{\linewidth}{@{}>{\raggedright\arraybackslash}Xrrr@{}}
\toprule
 & Closest study returned & Relevant study returned & Returned and cited \\
\midrule
Automatic document retrieval, keyword (up to 4) & 2/45 & 50/105 & 47/105 \\
Automatic document retrieval, embedding (10) & 18/45 & 59/105 & 55/105 \\
Automatic graph retrieval, keyword (up to 4) & 2/45 & 38/105 & 35/105 \\
Automatic graph retrieval, embedding (10) & 22/45 & 69/105 & 61/105 \\
\bottomrule
\end{tabularx}
\end{table}

\section{When memory helps, and when it hurts}
\label{sec:transfer}

Search changed whether the decision was correct in about a third of paired runs, for better and for worse about equally often; the three cases below show how both directions happen.

\subsection{The tally}

Each memory run can be paired with the no-memory run of the same model, task, and seed. Under the primary scoring rule there are 300 such pairs.

\begin{center}
\begin{tabular}{lr}
\toprule
Outcome & Pairs \\
\midrule
Memory fixed a wrong decision & 47 \\
Memory broke a correct decision & 46 \\
Decision unchanged, and correct & 136 \\
Decision unchanged, and wrong & 71 \\
\bottomrule
\end{tabular}
\end{center}

Memory changed whether the decision was correct in 93 of 300 pairs (31 percent), and it did so in the right and the wrong direction about equally often. DeepSeek is the exception, with 14 fixes and 4 breaks. Gemini fixed 10 and broke 12, Haiku 8 and 11, gpt-oss-20b 10 and 13, and Sonnet 5 and 6. Pairs share tasks, so the pairs are dependent.

I chose the next two cases, Case 1 and Case 2, because they are the pair where the same record led to opposite verdicts. I made that selection after seeing the results.

\subsection{Case 1: a stop the cited record leaves unexplained (RT-012)}
\label{sec:rt012}

\textbf{The idea.} The proposed study appeared to find that a teammate's drive before a pick-and-roll improved its efficiency. However, the finding depended on how pick-and-roll possessions were classified and did not hold under an alternative classification. The study was therefore classified as an artifact. The answer key says stop. The models could not see this result.

\textbf{The record.} An earlier study had identified problems with the play-type data, including unreliable classifications. This was a warning that the data might not be dependable, but it did not directly establish that the proposed finding was invalid. The numbers behind that warning are in Appendix~\ref{app:cases}.

\textbf{What happened.} Averaged across models, accuracy was 0.33 without memory, 0.67 with document search, and 0.60 with graph search. The audit was cited in 18 of the 30 memory runs, 14 of which were correct; 5 of the 12 memory runs that did not cite it were correct. Haiku was wrong in all three no-memory runs and correct in all six memory runs, and DeepSeek went from one correct no-memory run in three to six of six. The other three models did not show the pattern. Sonnet, correct in all three no-memory runs, was correct in only one of three document-search runs: it cited the audit in each, named label error as a risk, and chose to go ahead with a label audit written into the design rather than to stop (further details in Appendix~\ref{app:cases}).

\textbf{What it shows.} Citing the warning was associated with correct stops for some models. One model treated the same warning as a design step to fix, and the answer key scores that as wrong. The automatic relevance rule excludes this record from the pick-and-roll drive study (RT-012; Appendix~\ref{app:retrieval}).

\subsection{Case 2: the same record, the opposite verdict (RT-147)}
\label{sec:rt147}

\textbf{The idea.} A proposed analysis asks whether the effect on an opponent's next possession, after each kind of play, is large and stable enough to justify telling a team to avoid that play. The answer key says go ahead. The original study found that the type of play affected the opponent's next possession across all eight play types. However, the study also found a problem in how those results were used to recommend which plays a team should avoid. The underlying relationship was measurable, even though the recommendation system had a design problem.

\textbf{What happened.} Averaged across models, accuracy was 0.80 without memory, 0.47 with document search, and 0.60 with graph search. Under document search Haiku fell from 1.00 to 0.00, gpt-oss-20b from 1.00 to 0.33, and DeepSeek from 0.33 to 0.00; Gemini rose from 0.67 to 1.00, and Sonnet was correct in every run. The same labeling audit was cited in 11 memory runs on this task, and only 3 of them were correct; 13 of the 19 memory runs that did not cite it were correct. The prior study of the same question supports going ahead: it reports an opponent effect of $+0.12$ to $+0.14$ points per possession in all eight play types, replicated in the playoffs. Its identifier appeared in the retrieved text of four document-search runs, all correct, and in none of the runs of the three models whose accuracy dropped.

\textbf{What it shows.} Both tasks depend on how plays are classified, so the same warning applies to both, and nothing in it says why one study should stop and the other should go ahead. The stop on the pick-and-roll drive study (RT-012) is defensible because the apparent gain disappeared when the possessions were classified differently, and that result was outside the warning. A data-quality warning leaves the model to judge whether the problem invalidates a particular research question. The answer key separates the two tasks by results the models could see only later. A model applying the warning to both would stop both, and under the benchmark's rule one of those stops counts as wrong.

I keep the case as an example of how hard it is to judge whether a lesson transfers.

Case 3, the second-half improvement study (RT-179), where an analogous failure helped most models and was ignored by one, and the missing-player role study (RT-065) are in Appendix~\ref{app:cases}.

\section{Study 1b: supplying the relevant studies}
\label{sec:study1b}

Section~\ref{sec:embed} asks whether the search returned the earlier study. This section asks whether the decision gets better once those studies are already in hand. On the same ten tasks, each run was handed up to four relevant earlier studies, chosen by a rule committed before any run, through the same tool and settings as the main experiment (270 runs; design, selection rule, and one declared manual addition in Appendix~\ref{app:more}). A second version added a short scope annotation to each study.

\begin{table}[ht]
\centering
\caption{Accuracy score when the relevant studies are handed over. Document and graph are the keyword-search scores from the main experiment. Notes adds a short note on when each study's lesson applies. Sonnet has no notes runs. The last column is supplied studies minus no memory, with a 95\% uncertainty range.}
\label{tab:study1b}
\footnotesize
\setlength{\tabcolsep}{3pt}
\begin{tabularx}{\linewidth}{@{}l *{6}{>{\centering\arraybackslash}X}@{}}
\toprule
Model & \shortstack{No\\memory} & \shortstack{Document\\search} & \shortstack{Graph\\search} & Supplied & Notes & \shortstack{Supplied minus\\no memory} \\
\midrule
Gemini 2.5 Flash & 0.639 & 0.681 & 0.583 & 0.861 & 0.847 & \shortstack{$+0.222$\\{\scriptsize$[-0.056,+0.597]$}} \\
Claude Haiku 5.5 & 0.625 & 0.583 & 0.639 & 0.639 & 0.667 & \shortstack{$+0.014$\\{\scriptsize$[-0.250,+0.347]$}} \\
gpt-oss-20b & 0.542 & 0.528 & 0.528 & 0.625 & 0.750 & \shortstack{$+0.083$\\{\scriptsize$[-0.208,+0.431]$}} \\
Claude Sonnet 5.5 & 0.736 & 0.653 & 0.722 & 0.722 & n/a & \shortstack{$-0.014$\\{\scriptsize$[-0.125,+0.083]$}} \\
DeepSeek V4 Pro & 0.514 & 0.681 & 0.722 & 0.778 & 0.806 & \shortstack{$+0.264$\\{\scriptsize$[0.000,+0.556]$}} \\
\bottomrule
\end{tabularx}

\end{table}

\begin{table}[ht]
\centering
\caption{Correct rate on the three tasks named in the preregistered predictions, pooled over the five models.}
\label{tab:study1b_tasks}
\small
\begin{tabular}{lrrrr}
\toprule
Task & No memory & Document search & Graph search & Supplied \\
\midrule
RT-012 & 0.33 & 0.67 & 0.60 & 1.00 \\
RT-147 & 0.80 & 0.47 & 0.60 & 0.67 \\
RT-179 & 0.20 & 0.73 & 0.53 & 0.80 \\
\bottomrule
\end{tabular}
\end{table}

Supplying the studies raised the accuracy score for four of the five models (Table~\ref{tab:study1b}). Every uncertainty range includes no difference. That means no effect was detected. Supplied studies are above document search for every model and at or above graph search for every model. The notes on when a lesson applies do not change the result: the differences run from $-0.014$ to $+0.125$. On the three predicted tasks the direction matches the prediction (Table~\ref{tab:study1b_tasks}), which was written after reviewing Study 1's per-task results. The pick-and-roll drive study (RT-012) is the largest gain, with all 15 supplied runs correct, and which of its four studies produced the gain is unknown. Fifteen of 270 runs were unreadable, and the conclusions are unchanged under the other scoring rules.

\section{Discussion}
\label{sec:discussion}

\subsection{What the results show}

The measured findings are in Section~\ref{sec:results} and the cases. In short, a model can retrieve and accurately repeat an organization's research history without reliably knowing when that history applies to a new problem.

Three patterns plausibly contribute, though none was tested as a cause. First, retrieval often did not return the most relevant prior work. Second, a record that applied cleanly to one task was applied to another where the answer key still required a different call (Cases 1 and 2). Third, retrieved failures shifted models toward caution generally, helping on tasks that should be stopped and hurting on tasks that should go ahead (Section~\ref{sec:secondary}).

\subsection{Two points of failure, described}

The results suggest two possible reasons memory fell short. The experiment leaves open which was more important.

The first is finding the right studies. Keyword search returned the closest earlier study in only 2 of 45 runs per condition, and a returned study was nearly always cited (47 of 50 document-search runs and 35 of 38 graph-search runs). The required output asks every run to cite what it relied on. Better search left the accuracy score flat (Section~\ref{sec:embed}), which counts against weak search alone as the explanation.

The second is judging whether a study applies. Returning and citing a relevant study was not consistently linked to a correct decision. Among memory runs on tasks with a relevant study, runs that returned one were more accurate under document search (0.82 against 0.60) and less accurate under graph search (0.66 against 0.76). These comparisons mix models and tasks, so they show an association and cannot establish a cause. The follow-up experiment suggests that some models may have benefited because they were given the right studies directly, rather than having to find them. Other models showed little improvement, possibly because they were already performing relatively well or because they struggled to judge whether the studies applied. These are possible explanations.

\subsection{What this means for a team's research archive}

Sports organizations are building systems in which an AI agent reads the department's past analyses before a new project starts (Section~\ref{sec:related}). The implicit promise is that the agent will stop the team from re-running dead ends and flag flawed designs before months of analyst time are spent. This paper tests that promise on one analyst's archive, failures included.

The result is a caution about one design choice, not a measurement of any front office. On these ten decisions the agent repeated the archive accurately but could not reliably tell when a past failure applied, and that has costs in both directions. A false stop kills a promising idea because a similar study failed, and analyst time is the scarce resource. A false go-ahead misses a known data problem and sends weeks into a design that cannot work. Both errors arrive with accurate citations, so a wrong call looks as authoritative as a right one, and staff have little reason to doubt it.

Basketball analytics makes the problem sharper because many studies share one pipeline. A single finding about play-type labels, tracking coverage, or lineup definitions bears on dozens of later questions, and whether it invalidates a given question depends on how that question uses the data. The play-type labeling audit in Cases 1 and 2 is the example: it was a sound warning for one idea and misleading for another. A team's archive will contain many findings of this kind.

\subsection{Model choice}

Sonnet without memory (0.736) is 0.097 above Gemini and 0.194 above gpt-oss-20b without memory, and on these ten decisions the choice of model moved accuracy more than attaching this archive did. (Appendix~\ref{app:history} lists the provider and temperature differences between these models.) Gemini with the studies handed over scored 0.861, which mixes the model and the archive: it is Gemini plus the supplied studies. On these ten tasks, the gap between models is about as large as Gemini's gain from the supplied studies ($+0.222$) and larger than any gain from keyword search.

\section{Design implications for a team's archive}
\label{sec:practical}

These are hypotheses consistent with the results, written for a team building such a system. None was tested here.

\textbf{Record why a study stopped, and which questions its problem affects.} A stored study should say what it measured, and also why it failed and which downstream questions the same problem threatens. ``Play-type labels are unreliable'' does not tell an agent whether a new pick-and-roll question is affected. ``The apparent gain disappeared under an alternative classification of pick-and-roll possessions'' does.

\textbf{Treat the agent as a librarian, not the decision-maker.} In these runs the same agent retrieved a prior study and judged whether it applied, and it was right about as often as wrong. Having the agent surface prior work while an analyst decides whether it applies matches the revise step in case-based reasoning. Whether this improves decisions is untested.

\textbf{Keep the failures.} Many teams keep what worked and discard the rest. The archive here is a useful test because it kept untestable, artifact, and inconclusive studies, and those were the records the agent applied least reliably. A team that records only successes gives its agent nothing to learn from on the questions most likely to fail.

\textbf{Let the system say ``nothing relevant.''} Three of the ten tasks had no relevant earlier study, and the search tool returned its top results for any query. An analyst starting a genuinely new question will often have no close precedent, and an agent forced to retrieve something receives neighboring studies that may not apply. Shi et al.\ (2023) suggest such context can shift decisions.

\textbf{Check the retrieval rule against an analyst.} The automated relevance rule here missed a record that drove decisions (Case 1; Appendix~\ref{app:retrieval}). A team should compare any automatic measure of what is relevant against an analyst's judgment before relying on it.

\textbf{Expect lessons to expire.} Rules, tracking systems, and rosters change from season to season, so a lesson recorded under one setup may not hold under the next. This paper's tasks come from one four-season window and do not test expiry, but a record that states when its lesson applies would need to cover it.

\section{Limitations}
\label{sec:limits}

\textbf{Ten tasks and little room to improve.} Ten tasks cannot establish that a small difference is real, and the claim is that memory did not reliably help on these ten decisions. Only three tasks have a close earlier study, and two models score about 0.5 without memory, which is what a coin flip would give on a two-way decision, so the test may have little room to show a gain. A larger set of 60 to 100 tasks is the natural next step.

\textbf{What the answer key measures.} The key is the archive's own rule, and it records whether a study later turned out untestable or an artifact, which is known only after the task date. It rewards hindsight more than judgment: a sensible analyst could say stop on a study that later returned a solid ``no effect,'' and the key would score that as wrong. The missing-player role study (RT-065), where nearly every run in every condition disagreed with the key, shows the cost, and treating uncertain studies as go ahead is contestable, most clearly on the opponent next-possession study (RT-147, Case 2). Because both the answer key and the historical records use the archive's own outcome labels, a model might improve simply by recognizing and repeating those labels. The experiment leaves this possibility open. The task's own study is removed, but earlier studies of the same question stay in memory. The pick-and-roll drive study (RT-012) and the opponent next-possession study (RT-147), where the same record led to opposite decisions, argue against simply copying a label.

\textbf{The stop score.} As Appendix~\ref{app:scoring} explains, a run can be scored as a stop even when the model writes ``proceed with changes,'' so any shift toward caution moves the accuracy score even when the written recommendation does not change (Section~\ref{sec:secondary}). The models' estimates of whether a study could be tested ranked the tasks better with memory for four of five models (Section~\ref{sec:secondary}). An alternative would be to evaluate the models using their probability estimates rather than converting those estimates into a stop-or-go decision at a fixed 50\% cutoff. That might produce a different result, but the original scoring rule was established before the experiment and was retained. Results also depend on how unreadable outputs are scored, most for Sonnet (Appendix~\ref{app:scoring}).

\textbf{No human baseline.} No analyst blind to the outcomes judged these ten ideas, with or without the archive, so where the models sit relative to a reasonable analyst is unknown. Adding one such person would be a later comparison, and a useful one.

\textbf{Retrieval method.} A different embedding model, a hybrid, a reranker, or query reformulation might change the result (Section~\ref{sec:embed}). Handing the studies over directly uses the studies themselves.

\textbf{Coding and exclusion audits.} Mentioned, returned, and used (Appendix~\ref{app:retrieval}) were coded by deterministic rules over the saved run files. A check of the rules against hand coding is still open. A parsing error dropped every family tagged with a single risk code; the figures here use the correction (Appendix~\ref{app:retrieval}). At least one misclassification remains: the play-type labeling audit is excluded from the pick-and-roll drive study (RT-012) because it has no risk-code tag. An audit of the 450 primary runs, done after the results were known, confirmed that no retrieval returned a record from the task's own family (Appendix~\ref{app:snapshots}); because runs could only search the restricted snapshots, that verifies how the snapshots were built.

\textbf{Post-hoc analyses.} The tasks, the answer key, the scoring, and the run outputs were fixed before this paper was first written. Several analyses were computed after the results were known, from the saved run files: which studies were retrieved, whether any records were improperly included, the scoring-rule sensitivity analysis, the paired tally, and the citation analyses. Two comparisons were specified before their own runs: the embedding search, and the list of studies handed over in Study 1b. The predictions for Study 1b were written after reviewing the per-task results from the keyword conditions.

\textbf{Archive and domain.} The archive is heavy with failed and inconclusive work, which may strengthen the caution shift. A corpus drawn from a graveyard may push models toward caution, and transfer to fields whose records are cleaner is untested. All ten tasks are basketball research designs from one research program, so another sport or a live scouting workflow is untested.

\textbf{No accumulation.} Each run retrieves from a fixed archive. An unchecked loop that writes a wrong lesson back and retrieves it again is outside what this paper measures.

\section{Future directions}
\label{sec:future}

The following was left for later work. A relationship layer stored alongside the document archive could distinguish kinds of connection between investigations. Cases 1 and 2 indicate that a record would need to say whether a shared data problem actually affects the question being asked, and when its lesson applies. It would also need an independent check of each cited connection before a decision is accepted. Each would need more tasks than the present benchmark. For any follow-up, the definitions of mentioned, returned, and used should be fixed in the analysis code before the runs, and the relevance coding checked against hand coding on a sample, with agreement reported.

\section{Conclusion}
\label{sec:conclusion}

The central finding is that an AI model can retrieve and accurately repeat an organization's research history without reliably knowing when that history applies to a new problem. Across five models and 450 primary runs on ten basketball research ideas, memory changed whether the decision was correct in about a third of paired runs, fixing and breaking decisions about equally often, and did not improve the average. The same labeling audit was cited in runs that stopped one design and in runs that stopped another. The audit is a data-quality warning, and it does not say why one question should stop and the other should go ahead. Searching with embeddings returned relevant studies far more often and still did not help (Section~\ref{sec:embed}). Handing over the relevant studies pointed upward for four of five models, but ten tasks cannot confirm a difference of that size (Section~\ref{sec:study1b}). Whether a larger task set, a better retriever, or records that state when a lesson applies would change this is open, and a test with 60 to 100 tasks is the next step.

\section*{Availability}
\label{sec:availability}

The following are available at \url{https://github.com/ramizheman/research-harness-benchmark}.

\begin{itemize}\raggedright
\item the task-selection script, with seed 20261110 (\texttt{bench/\allowbreak tasks.py}, run by \texttt{scripts/\allowbreak 12\_generate\_replay\_tasks.py});
\item the two preregistrations for this paper: Study 1b (\texttt{study1b/\allowbreak prereg\_1b.md}) and the embedding search (\texttt{study1\_embed/\allowbreak prereg.md}). Study 1 has no separate preregistration file;
\item the 319-family archive (\texttt{memory\_v2/\allowbreak full/\allowbreak documents/\allowbreak families/}) and the per-task memory snapshots with their fact hashes (\texttt{memory\_v3/\allowbreak snapshots/}, one manifest per task);
\item the saved run files for the pilots, Study 1, Study 1b, and the embedding search (\texttt{runs/}), and the analysis scripts that reproduce every table (\texttt{study2/\allowbreak analysis.py}, \texttt{analysis\_fixes/\allowbreak recompute\_retrieval\_audit.py});
\item the scoring specification and the audit code, including the risk-code parsing correction (\texttt{bench/\allowbreak scoring.py}, \texttt{bench/\allowbreak retrieval\_audit.py});
\item the prompts used in the runs (\texttt{runner/\allowbreak agent.py}).
\end{itemize}

\section*{AI Disclosure}

Large language models (Claude, Anthropic; Grok, xAI) were used to help write the analysis code and draft this paper's text. The author directed the research question, supplied and validated the data, and reviewed and edited all results and text.

\begin{thebibliography}{99}

\bibitem{aamodt1994}
Aamodt, A., and Plaza, E. (1994).
Case-based reasoning: Foundational issues, methodological variations, and system approaches.
\textit{AI Communications}, 7(1), 39--59.
\url{https://doi.org/10.3233/AIC-1994-7104}

\bibitem{erhardt2016}
Erhardt, N., and Martin-Rios, C. (2016).
Knowledge management systems in sports: The role of organisational culture and structure.
\textit{Journal of Information \& Knowledge Management}, 15(2), 1650023.
\url{https://doi.org/10.1142/S0219649216500234}

\bibitem{football2024}
Science and Medicine in Football (2024).
Analytical capacity in professional and international football: a survey of 29 national federations and 32 clubs.
\url{https://doi.org/10.1080/24733938.2024.2341837}

\bibitem{lewis2020}
Lewis, P., Perez, E., Piktus, A., Petroni, F., Karpukhin, V., Goyal, N., K\"uttler, H., Lewis, M., Yih, W., Rockt\"aschel, T., Riedel, S., and Kiela, D. (2020).
Retrieval-augmented generation for knowledge-intensive NLP tasks.
\textit{Advances in Neural Information Processing Systems}, 33, 9459--9474.

\bibitem{mission2024}
Mission Cloud (2024).
Boston Celtics: cloud infrastructure case study.
\url{https://www.missioncloud.com/case-studies/boston-celtics}

\bibitem{sabres2026}
Buffalo Sabres (9 September 2026).
Buffalo Sabres and SprintAI announce an AI-powered hockey operations partnership.
\url{https://www.nhl.com/sabres/news/buffalo-sabres-sprint-ai-partnership}

\bibitem{shi2023}
Shi, F., Chen, X., Misra, K., Scales, N., Dohan, D., Chi, E. H., Sch\"arli, N., and Zhou, D. (2023).
Large language models can be easily distracted by irrelevant context.
\textit{Proceedings of the 40th International Conference on Machine Learning}, PMLR 202, 31210--31227.

\bibitem{ho2026}
Ho, S.-T., Liu, M., Nghiem, H., Tran, H., Agrawal, A., Lee, Y., Kang, M., Schoen, A. A., and Huang, F. (2026).
SoundnessBench: Can your AI scientist really tell good research ideas from bad ones?
\url{https://arxiv.org/abs/2605.30329}

\bibitem{tian2026}
Tian, Q., Yin, H., Xia, Y., Kong, Y., Liu, Z., Su, H., Chao, J., and Gu, X. (2026).
ForeSci: Evaluating LLM agents for forward-looking AI research judgment.
\url{https://arxiv.org/abs/2606.00644}

\bibitem{shinn2023}
Shinn, N., Cassano, F., Gopinath, A., Narasimhan, K., and Yao, S. (2023).
Reflexion: Language agents with verbal reinforcement learning.
\textit{Advances in Neural Information Processing Systems}, 36, 8634--8652.

\bibitem{sbj2025}
Sports Business Journal (18 August 2025).
Cleveland Cavaliers semantic search over scouting and statistical reports; San Antonio Spurs travel planning.
\url{https://www.sportsbusinessjournal.com}

\end{thebibliography}

\appendix

\section{Supplementary results and design history}
\label{app:more}

This appendix holds material moved out of the main text. In its tables, A is no memory, B is document search, C is graph search, O is supplied studies (the oracle condition of Study 1b), and OS is supplied studies with scope annotations; B-emb and C-emb are the embedding versions of B and C.

\subsection{Pilots, experiments, conditions, and protocol checks}
\label{app:history}

\subsubsection{Pilots and three experiments}

The protocol changed during development because early versions failed as procedures. Pilots 1 and 2 used one seed and were used to develop the protocol. Experiments 1 to 3 used the same protocol at seeds 1, 2, and 3, and form the 450 primary runs. The tasks, the answer key, and the scoring stayed the same across all of them. An audit of Pilot 2 found that none of the 160 items returned across its memory runs was a family record, so the intervention was not delivered. Experiment 1 fixed that before the primary runs.

The development record is at \url{https://github.com/ramizheman/research-harness-benchmark/blob/master/docs/paper1_supplement/pilots_development.md}.

\subsubsection{Three conditions}

Condition B ranks the full text of each family record by TF-IDF keyword matching against a query the model writes, and returns families that share at least one term with the query; shared terms are its only relevance test. Agents could make follow-up retrieval calls within a tool budget, so some runs saw more than four families. In Experiment 1, the first retrieval in B and C shared 1.4 of four families on average for the same task and seed; this overlap was logged only for Experiment 1.


\subsubsection{Models and protocol checks}

Gemini 2.5 Flash, Claude Haiku 5.5, and gpt-oss-20b were run at temperature 0.7. Claude Sonnet 5.5 and DeepSeek V4 Pro were run with extended thinking, which requires temperature 1.0. Sonnet and DeepSeek also differ from the other three models in provider, so a comparison across models mixes provider, temperature, and extended thinking. Each model's own temperature is fixed across conditions, so the within-model comparisons hold temperature fixed. Extended thinking was confirmed by the presence of reasoning traces passed between turns; the traces themselves were not analyzed.

The task text, answer key, scoring, and memory snapshots were shared. Before the live runs, unit tests checked that the document and the graph carried the same facts, that retrieval preceded any assessment, and that a citation had to match a record the tool had returned. One correction turn was allowed when a cited status did not match the record, and the pre-correction answer was kept.


\subsection{Per-model primary results}
\label{app:permodel}

\begin{table}[ht]
\centering
\caption{Balanced flag accuracy, primary scoring rule}
\label{tab:primary}
\small
\begin{tabular}{lrrrll}
\toprule
Model & A & B & C & B $-$ A [95\% CI], $p$ & C $-$ A [95\% CI], $p$ \\
\midrule
Gemini 2.5 Flash & 0.639 & 0.681 & 0.583 & $+0.042$ $[-0.264, +0.389]$, 0.91 & $-0.056$ $[-0.361, +0.292]$, 0.84 \\
Claude Haiku 5.5 & 0.625 & 0.583 & 0.639 & $-0.042$ $[-0.333, +0.319]$, 0.88 & $+0.014$ $[-0.333, +0.389]$, 1.00 \\
gpt-oss-20b & 0.542 & 0.528 & 0.528 & $-0.014$ $[-0.250, +0.250]$, 1.00 & $-0.014$ $[-0.333, +0.347]$, 1.00 \\
Claude Sonnet 5.5 & 0.736 & 0.653 & 0.722 & $-0.083$ $[-0.319, +0.111]$, 0.56 & $-0.014$ $[-0.125, +0.083]$, 1.00 \\
DeepSeek V4 Pro & 0.514 & 0.681 & 0.722 & $+0.167$ $[-0.097, +0.486]$, 0.44 & $+0.208$ $[+0.069, +0.389]$, 0.06 \\
\bottomrule
\end{tabular}

\vspace{0.4em}
\begin{minipage}{0.92\linewidth}\footnotesize
Intervals are percentile bootstrap intervals over the ten tasks (10,000 resamples); $p$ is the exact sign-flip $p$-value. After Holm correction across the ten contrasts, the smallest adjusted $p$-value is 0.62 (DeepSeek, C $-$ A).
\end{minipage}
\end{table}

\FloatBarrier

\subsection{What the models did with the record}

\subsubsection{Citation accuracy}

Every status checked against the record matched in final answers: Gemini 70 of 70 under document retrieval and 97 of 97 under graph retrieval, Haiku 86 of 86 and 98 of 98, gpt-oss-20b 77 of 77 and 80 of 80, Sonnet 96 of 96 and 103 of 103, and DeepSeek 94 of 94 and 97 of 97. Ten memory runs used the allowed correction turn. Condition A runs reported that no historical evidence was available.

One possible reading is that a stronger model reconstructs many of these failure modes without the archive, leaving less room for memory to help. The data leave that reading untested. DeepSeek's RT-012 runs are consistent with memory supplying a specific fact the model produced only with memory: it cited the labeling audit in four of six memory runs and was correct in all six. On these runs the pattern is too small to call.

A false prior claim is a memory run that says earlier work already answered substantially the same question on a task where no such work was in memory (eight tasks, 24 eligible runs per model and condition). Its rate was 25 to 50 percent of eligible runs under document retrieval and 42 to 67 percent under graph retrieval, and it was higher under graph retrieval for four of the five models.

\subsubsection{Caution, applied in both directions}

Table~\ref{tab:caution} shows Experiment 1. Retrieved history lowered Gemini's stated probability that the proposed study was testable, both on tasks that should be flagged and on tasks that should not.

\begin{table}[ht]
\centering
\caption{Caution shift, Gemini 2.5 Flash, Experiment 1 (primary scoring rule, 30 runs per condition)}
\label{tab:caution}
\begin{tabular}{lccc}
\toprule
 & A & B & C \\
\midrule
Mean $p$(testable), should flag & 0.48 & 0.26 & 0.40 \\
Mean $p$(testable), should not flag & 0.61 & 0.51 & 0.51 \\
Flag rate, should flag & 0.50 & 0.75 & 0.50 \\
Flag rate, should not flag & 0.22 & 0.39 & 0.33 \\
\bottomrule
\end{tabular}
\end{table}

\FloatBarrier

\subsection{Full case accounts}
\label{app:cases}

\subsubsection{Additional details, Case 1}

\textbf{What the archive contained.} An earlier labeling audit, FAM:SF-CHT-025, was in memory. Its record states that foul-ending possessions inflated the efficiency of the OTHER category (36 percent of those possessions ended in fouls), that drives sat 12 points below the foul benchmark, that pick-and-rolls showed a 22 percent foul rate, and that only 335 of 50,668 possessions (0.7 percent) mapped to a real play type through the initiation-type field. It records the outcome as an artifact. For RT-012, the document-retrieval result in many runs included FAM:SF-CHT-025.

\textbf{What the models did.} The clean pattern, mostly wrong without memory and correct in every memory run, occurred for two models in four of the ten model-condition cells. Haiku was wrong in all three no-memory runs and correct in all six memory runs, citing FAM:SF-CHT-025 in each; in its document-retrieval runs it named the 0.7 percent coverage as a blocking constraint. Gemini was correct in one of six memory runs. gpt-oss-20b was correct in three of six.

\textbf{What the case shows.} The automated relevance rule in Appendix~\ref{app:retrieval} excludes FAM:SF-CHT-025 from RT-012, because the record carries no risk-code tag; a human reader would plausibly treat it as a direct dependency. The appendix's retrieval rates for this task should be read with that in mind.

\subsubsection{Additional details, Case 2}

\textbf{What the models did.} Haiku also fell to 0.00 under graph retrieval, and gpt-oss-20b to 0.67.

{\sloppy Haiku's three document-retrieval runs all cited FAM:SF-CHT-025, the outcome-selected suppression-label artifact FAM:SF-CHT-086, and the inferred defensive-scheme labels recorded as untestable in FAM:SF-APPA-005.\par} It added that the play-type labels had failed audits.

The prior study of the same question, SF07, was in memory. Its identifier appeared in the retrieved text of four document-retrieval runs (two from Gemini and two from Sonnet), all of them correct. Gemini cited it in one document-retrieval run and two graph-retrieval runs.

\subsubsection{Case 3: an analogous failure, used by most models and ignored by one (RT-179)}

\textbf{Question.} A proposed analysis asks whether a team's ability to limit opponent scoring improvement across halves is a stable team characteristic. The answer key says stop. When the data was divided into two halves, the measure showed virtually no consistency between them ($r = -0.001$). In other words, it could not reliably measure the team characteristic the study was trying to identify.

\textbf{What the models did.} Accuracy rose from 0.20 without memory to 0.73 with document retrieval and 0.53 with graph retrieval. Gemini, gpt-oss-20b, and DeepSeek were wrong in every no-memory run, and each was correct in at least two of three document-retrieval runs. Sonnet was wrong in every run in every condition, including graph-retrieval runs in which it returned and cited families the audit counts as relevant.

\textbf{What the archive contained.} No prior family studies this question, but seven lesson-tier families do: prior untestable or uncertain studies tagged with a measurement-reliability or mechanical-floor risk. Every memory run on this task, under both conditions, returned at least one of them. In 14 of 15 document-retrieval runs, the models cited FAM:SF-APPS-085, a study that found a different trait (a player's teammate-path quality) too unreliable to measure, with a split-half correlation of 0.215. Eleven of those 14 runs were correct. In graph retrieval the same family was cited in 10 runs, 6 of them correct, and accuracy was lower than under document retrieval even though a lesson-tier family was returned in every run.

\textbf{What the case shows.} On this task the analogy appears to have helped: a reliability failure for a different trait was applied to a new trait with the same measurement problem, and most models stopped correctly. Two cautions apply. A general caution shift of the kind in Table~\ref{tab:caution} would also raise accuracy on a task that should be flagged, and the run data leave the two explanations mixed. And one model, Sonnet, returned and cited the same families and still proceeded in every run.

\subsubsection{Tasks with no relevant prior}

RT-065 asks whether pre-absence profiles identify which teammate absorbs a departing player's role. The answer key says go ahead (an adequately powered test rejected the hypothesis), but accuracy was 0.07 without memory, 0.07 with documents, and 0.13 with the graph. The archive holds no relevant family for this question in either tier. Because the error appears without memory, retrieval did not cause it, and memory did not correct it.


\subsection{Study 1b: oracle retrieval}

Design details for Section~\ref{sec:study1b}. The gold families, annotations, and analysis were fixed in a preregistration committed before any run.

\subsubsection{Design}

The supplied-studies conditions were added alongside no memory, document search, and graph search. In condition O (oracle), each run received a frozen list of gold relevant prior studies (Section~\ref{sec:tasks}) for its task through the same retrieval tool, protocol, prompts, and settings as Study 1; the model could also read any of those families in full. Each list held up to four families: the first two question-tier and the first two lesson-tier families under the corrected audit of Appendix~\ref{app:retrieval}, by a rule committed before the list was computed. One declared manual addition was made: FAM:SF-CHT-025 was added to the lists for RT-012 (where it replaced FAM:SF02) and RT-147, because the automated rule cannot select it (Appendix~\ref{app:retrieval}). The lists were not changed after any run. Seven tasks have gold families (two to four each). The three tasks with none (RT-065, RT-101, RT-146) serve as a noise reference, because O and A are the same condition there apart from the extra read tool. Documents in O use the condition B format, so O$-$C mixes the effect of which families are retrieved with the effect of format.

In condition OS, each gold family also carried a short scope annotation stating which data, estimands, and effect sizes its recorded defect does and does not threaten. There are 23 annotations, written only from each family's own record, with identical text wherever a family appears. They were checked mechanically for shared phrases with the answer key and with the task-own records. The annotator knew the answer key and the Study 1 results, and the ``applies to'' sentences are the annotator's inference from the records, so OS is exploratory. OS also carries 729 to 1,784 more characters in the first retrieval than O, so OS$-$O measures annotation plus extra length. OS was run for four models; Sonnet 5.5 was run on O only. Each condition used three seeds per task and model (270 runs).

\subsubsection{Results}

The Holm-adjusted $p$-values for O$-$A were 0.94 or higher.

Against Study 1's own search the direction is more uniform. O exceeded B for all five models ($+0.056$ to $+0.181$) and was at or above C for all five ($0.000$ to $+0.278$). All of these intervals include zero, and three lower bounds sit at $0.000$ (O$-$B for Haiku, O$-$C for Gemini, and under the valid-only rule O$-$A for Gemini).

The predictions, written after reviewing Study 1's per-task results, were that O would move RT-012 and RT-179 toward the answer key and RT-147 away from it, and that annotations would not clearly change any of the three. On the four models that ran OS, the change from O to OS was $0.00$ for RT-012, $0.00$ for RT-179, and $+0.17$ for RT-147; the preregistration set no numeric threshold for these calls, so they are directional. All 12 OS runs on RT-012 were correct.

Scope annotations did not change decisions by any preregistered criterion. No model shifted its flag rate by two thirds or more on two or more tasks. On the three tasks without gold families, where O and OS are identical, flag shares changed by up to 0.67 for gpt-oss and 0.33 for each other model. On annotated tasks the largest changes were 1.0 (gpt-oss, RT-112) and 0.67 (Haiku, RT-147). The one task where annotations moved the pooled rate is RT-147 (O 0.58, OS 0.75 on the four OS models), which is within the range that run-to-run noise produces here. Because the annotation for FAM:SF-CHT-025 is the same text for RT-012 and RT-147 and holds no information about effect size, this experiment does not test the scope-condition idea in the form Study 1 suggested.

\subsubsection{Limits and deviations}

Two of the 270 runs had no scorable flag. On the three no-gold tasks, O and A differ by up to 0.67 in a single cell (for example DeepSeek on RT-146), a reminder of how noisy cells of three runs are. The predictions were written after reviewing Study 1's per-task results and revised once before any Study 1b run. All analyses were computed from saved runs by a script that uses the same methods as Study 1, and the Study 1 columns reproduce the published table.


The per-model values are in Table~\ref{tab:study1b}, and the three predicted tasks are in Table~\ref{tab:study1b_tasks}.

\subsection{Embedding search: design and limits}
\label{app:embed}

\subsubsection{What the embedding conditions were}

Embedding search ranks each study by the cosine similarity of bge-small-en-v1.5 embeddings, with no relevance cutoff. Each search returns ten studies. The tool-result limit is 14,000 characters so that ten studies fit. In graph search the same embeddings are used for concept nodes and for each study's question and hypothesis. Keyword search returns up to four studies, and its tool-result limit is 6,000 characters. There were 300 embedding runs.

\subsubsection{Results}

Retrieval improved a great deal on the closest study and moderately on relevant studies overall (Table~\ref{tab:embed_retr}). The gain on a relevant prior study is 9 runs for document search and 31 for graph search.

Accuracy did not improve (Tables~\ref{tab:embed_acc} and~\ref{tab:embed_diff}). Under the primary scoring rule the ten embedding-minus-no-memory differences range from $-0.125$ to $+0.125$, the smallest unadjusted sign-flip $p$ is 0.25 (Sonnet, document search), and every Holm-adjusted $p$ is 1.00. Pooled over models the differences are $-0.053$ (95\% interval $-0.256$ to $0.181$) for document search and $-0.033$ ($-0.228$ to $0.181$) for graph search. Embedding accuracy was also lower than keyword accuracy for the same condition in all ten model-condition comparisons (by 0.014 to 0.125), and pooled the difference is $-0.067$ ($-0.158$ to $0.008$) for document search and $-0.061$ ($-0.197$ to $0.075$) for graph search, neither detectable after correction. Under the strict scoring rule, pooled embedding document search was 0.089 below keyword document search (interval $-0.153$ to $-0.033$, unadjusted $p = 0.016$, Holm-adjusted $0.062$), which does not meet the 0.05 threshold. Sixteen of 150 embedding document-search runs and 6 of 150 embedding graph-search runs were invalid, 8 of the 16 for Haiku, and one graph-search run had no scorable flag.

\begin{table}[ht]
\centering
\caption{Balanced accuracy under the primary scoring rule. A is no memory, B is document search, and C is graph search. B-emb and C-emb are the embedding versions.}
\label{tab:embed_acc}
\small
\begin{tabular}{lrrrrr}
\toprule
Model & A & B & B-emb & C & C-emb \\
\midrule
Gemini & 0.639 & 0.681 & 0.653 & 0.583 & 0.569 \\
Haiku & 0.625 & 0.583 & 0.500 & 0.639 & 0.542 \\
gpt-oss & 0.542 & 0.528 & 0.444 & 0.528 & 0.431 \\
Sonnet & 0.736 & 0.653 & 0.639 & 0.722 & 0.708 \\
DeepSeek & 0.514 & 0.681 & 0.556 & 0.722 & 0.639 \\
\bottomrule
\end{tabular}
\end{table}

\begin{table}[ht]
\centering
\caption{Embedding minus no memory (A), with 95\% bootstrap intervals over tasks and Holm-adjusted sign-flip $p$ across the ten contrasts.}
\label{tab:embed_diff}
\small
\begin{tabular}{lll}
\toprule
Model & B-emb $-$ A [95\% interval], $p_{\mathrm{Holm}}$ & C-emb $-$ A [95\% interval], $p_{\mathrm{Holm}}$ \\
\midrule
Gemini & $+$0.014 [$-$0.361, 0.458], 1.00 & $-$0.069 [$-$0.375, 0.292], 1.00 \\
Haiku & $-$0.125 [$-$0.445, 0.264], 1.00 & $-$0.083 [$-$0.375, 0.278], 1.00 \\
gpt-oss & $-$0.097 [$-$0.361, 0.264], 1.00 & $-$0.111 [$-$0.347, 0.111], 1.00 \\
Sonnet & $-$0.097 [$-$0.194, 0.000], 1.00 & $-$0.028 [$-$0.167, 0.083], 1.00 \\
DeepSeek & $+$0.042 [$-$0.194, 0.361], 1.00 & $+$0.125 [$-$0.139, 0.458], 1.00 \\
\bottomrule
\end{tabular}
\end{table}

Finding more relevant studies did not improve decisions. In this design, finding the relevant studies far more often left the decisions flat, which counts against weak retrieval alone as the explanation for the flat result. The comparison uses one embedding model and one cutoff. A different embedding model, a hybrid, a reranker, or an agent that reformulates its queries might still do better (Section~\ref{sec:embed}). It is also consistent with the applicability problem described in Section~\ref{sec:discussion}, and this experiment leaves that untested.

\FloatBarrier







\subsection{Task selection}
\label{app:selection}

The answer key is in Section~\ref{sec:tasks}. The ten tasks were chosen by a script. It drew from 196 eligible study families, first building a balanced main set of 78 (16 supported, 16 rejected, 16 untestable, 16 artifact, 14 uncertain), then a ten-task set with two per outcome group. Candidates were ranked by the archive's confidence score, with bonuses for a self-contained question and a reliable date. A seeded coin flip (seed 20261110) broke near ties, one task per question cluster was kept, and repeats of a topic were capped. Pilots 1 and 2 ran on these ten tasks before the 450 primary runs. I read the pilot outputs while developing the protocol and kept the task list afterward.

\subsection{Task details}

\begin{table}[ht]
\centering
\caption{The ten benchmark tasks: research questions, expected decisions, and recorded results}
\label{tab:tasks}
\small
\begin{tabular}{p{0.6cm}p{1.6cm}p{4.4cm}p{1.0cm}p{5.2cm}}
\toprule
Task & Topic & Research question (as posed) & Decision & Recorded result (not shown to the models) \\
\midrule
RT-012 & Play-type labeling & Does a teammate drive preceding a pick-and-roll raise that PNR's efficiency? & Stop & Artifact. The apparent gain did not hold under an alternative classification of pick-and-roll possessions. \\
\addlinespace
RT-014 & Defensive suppression & Does adhering to a suppression-triggered recommended play-type mix predict game outcomes? & Stop & Artifact. The entire apparent win gain ($+11.5$ pp) was second-half PPP leakage from the same game. Ex-ante effect $-2.7\%$, $p = 0.47$. \\
\addlinespace
RT-065 & Player absence & Can pre-absence initiation profiles, role similarity, or minutes identify which teammate will absorb a departing player's offensive role? & Go ahead & Rejected. 148 absence spells. Profile match 19\% (chance 14\%); 0 of 6 high-load cases correct. \\
\addlinespace
RT-066 & Player absence & Does the degree of offensive reorganization during an absence predict the change in team scoring efficiency? & Go ahead & Rejected. No relationship in 986 player-seasons; primary analysis $\rho = -0.040$, 95\% CI $[-0.21, +0.17]$, $p = 0.62$. \\
\addlinespace
RT-101 & Structural roles & Does lineup role coverage vary within teams in a way that is separable from raw talent? & Go ahead & Supported. 72 team-seasons; coverage--talent $|\rho| = 0.042$. All three gates passed. \\
\addlinespace
RT-112 & Lineup effects & Does changing the offensive five-man lineup detectably alter shot-type mix? & Go ahead & Supported. Own-lineup three-point rate $+3.1$ pp (SE 1.2); opponent-only changes null. \\
\addlinespace
RT-146 & Defensive suppression & After pick-and-roll efficiency is suppressed, does pull-up shooting provide a reliably positive counter in at least 70\% of cases? & Go ahead & Uncertain. Zero branches met the rule (at least 70\% of units, median lift above 2 pp). \\
\addlinespace
RT-147 & Play-type labeling & Are opponent next-possession effects after each play type large and stable enough to justify the AVOID labels the system generates? & Go ahead & Uncertain. Effect about $+0.12$ PPP in 8 of 8 play types; using it in the label score put PNR in AVOID even at 0.93--0.95 own PPP. Recorded mechanism: artifact. \\
\addlinespace
RT-179 & Defensive suppression & Is a team's ability to limit opponent scoring improvement across halves a reliable team characteristic? & Stop & Untestable. Split-half $r = -0.001$; a detectable effect needs $r \approx 0.18$. \\
\addlinespace
RT-183 & Structural roles & Do cut-heavy role players maintain their role more consistently after changing teams than drive-heavy players? & Stop & Untestable. $n = 24$, $p = 0.13$. \\
\bottomrule
\end{tabular}
\end{table}

\vspace{0.4em}
\begin{minipage}{0.92\linewidth}\footnotesize
Terms used in the table above. PNR is pick-and-roll, a ball handler using a teammate's screen. PPP is points per possession. AVOID is a label the recommendation system attaches to a play type it advises against.
\end{minipage}

\FloatBarrier

\section{Task-level flag accuracy}
\label{app:tasks}

Table~\ref{tab:cells} reports the share of seeds with a correct flag decision under the primary scoring rule. Each cell is the mean of three seeds, and the balanced accuracies in Table~\ref{tab:primary} are computed from these cells.

\begin{table}[ht]
\centering
\caption{Correct-decision rate by task, condition, and model (primary scoring rule)}
\label{tab:cells}
\footnotesize
\begin{tabular}{llccccc}
\toprule
Task & Stop & Gemini & Haiku & gpt-oss & Sonnet & DeepSeek \\
\midrule
RT-012 & Yes & 0 / .33 / 0 & 0 / 1 / 1 & .33 / .67 / .33 & 1 / .33 / .67 & .33 / 1 / 1 \\
RT-014 & Yes & 1 / 1 / .67 & 1 / 1 / 1 & 1 / 1 / 1 & 1 / .67 / 1 & 1 / .67 / 1 \\
RT-065 & No & 0 / 0 / .33 & 0 / 0 / 0 & .33 / 0 / .33 & 0 / .33 / 0 & 0 / 0 / 0 \\
RT-066 & No & 1 / .67 / 1 & 0 / 0 / .67 & 0 / 0 / .33 & .67 / 1 / .67 & .67 / .67 / 1 \\
RT-101 & No & 1 / 0 / .33 & 0 / 0 / 0 & .33 / 0 / 0 & 1 / 1 / 1 & 1 / 1 / 1 \\
RT-112 & No & 1 / 1 / 1 & 1 / .33 / 1 & .67 / 1 / 1 & 1 / 1 / 1 & .67 / 1 / .67 \\
RT-146 & No & 1 / 1 / .67 & 1 / .67 / 1 & .67 / 0 / 0 & .67 / 1 / 1 & 0 / 0 / .33 \\
RT-147 & No & .67 / 1 / .67 & 1 / 0 / 0 & 1 / .33 / .67 & 1 / 1 / 1 & .33 / 0 / .67 \\
RT-179 & Yes & 0 / 1 / 1 & 1 / 1 / .33 & 0 / .67 / 1 & 0 / 0 / 0 & 0 / 1 / .33 \\
RT-183 & Yes & 1 / .67 / .33 & 1 / 1 / 1 & 1 / 1 / .33 & 1 / .67 / 1 & 1 / 1 / 1 \\
\bottomrule
\end{tabular}

\vspace{0.4em}
\begin{minipage}{0.92\linewidth}\footnotesize
Each entry is condition A / B / C. Stop is whether the correct decision is to stop the design. Source: the scored run files for Experiments 1 to 3.
\end{minipage}
\end{table}

\FloatBarrier

\section{Retrieval and use of relevant families}
\label{app:retrieval}

\subsection*{Definitions}

All measures are computed by deterministic rules from the saved run files.

\textbf{Question-tier relevant family.} A family in the task's snapshot that is in the same research-question cluster as the task, is named in the answer key as a knowable precedent or a prior study of the same question, or has text similarity to the task above the archive's merge threshold.

\textbf{Lesson-tier relevant family.} A family in the task's snapshot recorded as untestable, artifact, or uncertain that carries one of the task's expected risk codes. RT-147 has no expected risk codes, so its lesson tier is empty by construction.

\textbf{Mentioned.} The study's identifier appears anywhere in the text the search returned, including inside another study.

\textbf{Returned.} The search actually returned that study.

\textbf{Used.} The study was returned, the model cited it, and the explanation is at least 30 characters long. This is only a basic indication that the model referred to the study. The required output asks for citations, so nearly any returned study that is cited passes.

Per-task retrieval shares are at \url{https://github.com/ramizheman/research-harness-benchmark/blob/master/docs/paper1_supplement/table16_retrieval.md}, with the same values in \url{https://github.com/ramizheman/research-harness-benchmark/blob/master/docs/paper1_supplement/table16_retrieval.csv}.

\subsection*{Summary}

Pooled over the seven tasks that have at least one relevant family (105 runs per condition), a relevant family was mentioned in 65 document-retrieval runs and 47 graph-retrieval runs, returned in 50 and 38, and used in 47 and 35. Used as a share of returned is 0.94 for documents and 0.92 for the graph. Used as a share of mentioned is 0.72 and 0.74; the gap between mentioned and returned reflects identifiers that appear in the text of other families without the family itself being returned.

Question-tier families were rarely returned. RT-012's question-tier family was never returned on that task. SF07 was returned in one document-retrieval run and two graph-retrieval runs on RT-147, all from Gemini; its identifier was mentioned in four document-retrieval runs, two from Gemini and two from Sonnet. SF-CHT-057 was returned in one document-retrieval run on RT-183, from DeepSeek, and in no graph-retrieval run. Across the three tasks with a question-tier family, that family was therefore returned in 2 of 45 runs under each method.

Mentioning, returning, or using a relevant family was not consistently associated with accuracy. Under the primary rule, runs that returned a relevant family were correct in 0.82 of document-retrieval runs and 0.66 of graph-retrieval runs, against 0.60 and 0.76 for runs that did not.

\subsection*{Correction to the audit code}

The audit code that produced the saved per-run files contained a parsing error. In the snapshot graph, a family's risk codes are stored as a list when there are several and as a plain string when there is one. The parser accepted only lists, so every family with a single risk code had no risk codes for the audit, and none could enter the lesson tier. In the RT-179 snapshot, 106 of the 151 families with risk codes store them as a single string, including FAM:SF-APPS-085 (MEASURE\_RELIABILITY). All figures in this appendix and in the main text are recomputed from the saved raw runs with only that parsing corrected. Before applying the correction, the recomputation was checked against the saved audit: with the original parser it reproduces all 300 memory-run rows exactly. The correction enlarges the lesson tier for six tasks (RT-012 from 3 to 5 families, RT-014 from 6 to 8, RT-066 from 21 to 49, RT-112 from 7 to 9, RT-179 from 6 to 7, RT-183 from 21 to 43), raises the returned rate on RT-012, RT-066, RT-179, and RT-183, and removes an apparent reversal in which document retrieval seemed to return nothing relevant for RT-179. The audit reads family attributes from the v2 rendering of each task's permitted records, which carries the stratum and risk fields; the v3 rendering the agents searched stores stratum and risk under different field names, so the audit uses the v2 rendering, where those fields keep their names.

\subsection*{Known limits of the relevance rule}

FAM:SF-CHT-025, the labeling audit cited in most RT-012 memory runs, is excluded from RT-012 because it carries no risk-code tag. CHT-039, a series-level study of AVOID share and winning, is excluded from RT-147 under either tier. The rule has not been checked against hand coding.

One difference between conditions remains: document retrieval returned a relevant family for RT-112 in every run, and graph retrieval in 2 of 15. It is reported as an observation from one task.

\FloatBarrier

\section{Memory snapshots}
\label{app:snapshots}

\begin{table}[ht]
\centering
\caption{Per-task memory snapshots}
\label{tab:snapshots}
\small
\begin{tabular}{llrrl}
\toprule
Task & Task date & Records & Families & Own family (excluded) \\
\midrule
RT-012 & 2026-06-14 & 166 & 148 & SF-CHT-032 \\
RT-014 & 2026-06-16 & 187 & 168 & SF-CHT-045 \\
RT-065 & 2026-11-01 & 317 & 266 & SF-PAP-056 \\
RT-066 & 2026-11-01 & 317 & 266 & SF-PAP-058 \\
RT-101 & 2026-07-13 & 238 & 203 & SF-PAP-009 \\
RT-112 & 2026-11-02 & 349 & 289 & SF-PAP-080 \\
RT-146 & 2026-05-23 & 104 & 96 & SF-CHT-026 \\
RT-147 & 2026-06-19 & 194 & 175 & SF-CHT-044 \\
RT-179 & 2026-10-06 & 293 & 248 & SF-CHT-064 \\
RT-183 & 2026-08-08 & 264 & 226 & SF-PAP-021 \\
\bottomrule
\end{tabular}

\vspace{0.4em}
\begin{minipage}{0.92\linewidth}\footnotesize
Records and Families are the permitted pool each task's runs could search. The task's own family, records citing it, probable duplicates, and records sharing three or more distinctive result numbers were also removed. Each snapshot has a manifest with a fact hash, which every run logged.
\end{minipage}
\end{table}


\FloatBarrier

\section{Scoring details, sensitivity, and invalid outputs}
\label{app:scoring}

This appendix collects the scoring mechanics, the sensitivity checks, and the related counts.

\subsection*{How a recommendation becomes a scored decision}

These are two different things. The model writes a recommendation. The benchmark then scores that recommendation as a stop or a go-ahead, and the two can disagree.

The model chooses one of three written recommendations: proceed, proceed with changes, or do not proceed. In that format, proceed means go ahead and do not proceed means stop.

The benchmark scores the run as a stop if any one of these is true:

\begin{itemize}
\item the written recommendation is ``do not proceed'';
\item the model says the study has less than a 50 percent chance of being testable;
\item the model assigns more than a 50 percent probability to the study being untestable or an artifact.
\end{itemize}

Otherwise the run is scored as a go-ahead. A model can write ``proceed with changes'' and still be scored as a stop, if either probability meets its threshold. The written recommendation and the scored decision can therefore contradict each other. This rule was fixed before the runs. It is why a general shift toward caution changes accuracy even when the written word stays ``proceed.''

The answer key is in Section~\ref{sec:tasks}.

Balanced accuracy is the average of two numbers: accuracy on the four stop tasks, and accuracy on the six go-ahead tasks. Averaging them keeps the six go-ahead tasks from dominating the score. Seeds are repeated runs of the same task. The ten tasks are the unit of analysis.

The primary rule scores invalid outputs where a decision can still be read. An output that fails schema validation but still yields a stop-or-go decision is scored on that decision, and an output from which no decision can be read counts as wrong. Of 32 invalid outputs in the primary runs, 31 yielded a decision; one (gpt-oss-20b, graph search) did not. Two sensitivity rules are reported alongside: strict, in which every invalid output counts as wrong, and valid-only, in which invalid outputs are dropped.

I tested whether the observed differences could reasonably have occurred by chance. Because ten comparisons were made, the tests were also adjusted to reduce the risk of finding an apparent improvement simply by making many comparisons. The exact procedure is a sign-flip test over all $2^{10}$ assignments of which tasks count for or against the difference, with Holm's adjustment across the ten comparisons. The ranges are percentile bootstrap intervals from 10,000 resamples of the ten task-level differences. With ten tasks they are about 0.6 wide, and a difference smaller than roughly 0.3 falls inside that width.

\subsection*{Sensitivity to the scoring rule}

Table~\ref{tab:sensitivity} reports the same contrasts under the two sensitivity rules. The rules matter most for Sonnet, which produced six invalid outputs under document retrieval and none without memory. Under the strict rule, Sonnet's document-retrieval difference is $-0.194$ with an interval that excludes zero ($[-0.389, -0.014]$, $p = 0.12$). Under the valid-only rule it is zero. DeepSeek's graph-retrieval interval excludes zero under the valid-only rule and not under the strict rule. No contrast has $p < 0.05$ under any rule.

\begin{table}[ht]
\centering
\caption{Memory contrasts under the two sensitivity rules}
\label{tab:sensitivity}
\small
\begin{tabular}{lllll}
\toprule
 & \multicolumn{2}{c}{Strict (invalid = wrong)} & \multicolumn{2}{c}{Valid outputs only} \\
\cmidrule(lr){2-3}\cmidrule(lr){4-5}
Model & B $-$ A & C $-$ A & B $-$ A & C $-$ A \\
\midrule
Gemini & $+0.000$ $[-0.264, +0.264]$ & $-0.083$ $[-0.389, +0.278]$ & $+0.042$ $[-0.264, +0.389]$ & $-0.069$ $[-0.375, +0.292]$ \\
Haiku & $-0.042$ $[-0.333, +0.319]$ & $-0.014$ $[-0.361, +0.375]$ & $-0.056$ $[-0.347, +0.306]$ & $+0.035$ $[-0.292, +0.403]$ \\
gpt-oss-20b & $+0.111$ $[-0.111, +0.347]$ & $+0.056$ $[-0.278, +0.417]$ & $+0.042$ $[-0.194, +0.292]$ & $+0.125$ $[-0.111, +0.417]$ \\
Sonnet & $-0.194$ $[-0.389, -0.014]$ & $-0.069$ $[-0.167, +0.000]$ & $+0.000$ $[-0.222, +0.167]$ & $+0.014$ $[-0.042, +0.083]$ \\
DeepSeek & $+0.139$ $[-0.125, +0.458]$ & $+0.167$ $[-0.028, +0.361]$ & $+0.167$ $[-0.097, +0.486]$ & $+0.208$ $[+0.069, +0.389]$ \\
\bottomrule
\end{tabular}

\vspace{0.4em}
\begin{minipage}{0.92\linewidth}\footnotesize
Smallest sign-flip $p$-values: strict, Sonnet B $-$ A, 0.12; valid only, DeepSeek C $-$ A, 0.06.
\end{minipage}
\end{table}

\FloatBarrier

\end{document}
